On-policy distillation (OPD) has rapidly emerged as a core technique for large language model (LLM) post-training \citep{lu2025onpolicydistillation}. OPD allows a smaller student model to align its policy with a stronger teacher model by learning directly from on-policy rollouts, enabling the student to effectively inherit complex reasoning capabilities. Consequently, recent pioneering industry efforts, including Qwen3 \citep{yang2025qwen3}, MiMo \citep{xiao2026mimo}, and GLM-5 \citep{zeng2026glm}, have integrated OPD into their post-training pipelines, establishing it as an indispensable complement to supervised fine-tuning (SFT) and reinforcement learning with verified reward (RLVR).

However, while researchers have focused heavily on designing better training algorithms (e.g. EOPD \citep{jin2026entropy} and REOPOLD \citep{ko2026scaling}), the data side of OPD remains relatively underexplored. Specifically, How much data is truly necessary? What data is most effective? And what actually drives the student model's improvement in the training data? Answering these questions is of significant practical value because in many real-world scenarios, high-quality data are extremely scarce or expensive to collect, particularly in highly specialized domains such as medicine and engineering. This bottleneck makes a systematic study of data necessity highly important. 

To answer these questions, we first investigate a special setting where we train OPD on only one example, namely \textbf{1-shot OPD}. Our empirical evaluations demonstrate that 1-shot OPD is consistently effective across all sampled training examples, which is also observed in a concurrent work \citep{fu2026rethinking}. Furthermore, we show that its training is exceptionally stable and the validation performance remains remarkably stable up to even 2,000 steps without overfitting or catastrophic policy collapse. We further analyze why 1-shot OPD works so well and find that the student policy aligns closest with the teacher on structural reasoning tokens (e.g. ``Alternative'', ``Wait'') after distillation. This indicates that the student can successfully learn and generalize the teacher’s reasoning patterns by distilling on only a single training example.

Then, we investigate what kind of data makes OPD most effective. We discover a clear trend: training on harder examples often yields better performance than training on easy or medium ones. Interestingly, even when these examples are so difficult that they completely exceed the capabilities of both the teacher and student models, keeping the training accuracy strictly at 0\%, training on these harder examples still allows the validation accuracy to continuously improve. We further analyze why the example difficulty can influence the OPD and what actually drives the student model's improvement in the training data. We find that this improvement is not driven by high-entropy tokens, but is instead determined by the longer CoT. Specifically, we observe that training with longer CoTs helps the student maintain a closer alignment with the teacher across long horizons. Moreover, these longer paths help the model acquire thinking patterns that are rarely learned from short CoTs, such as self-reflection (e.g., "Alternatively").

Based on these findings, we propose a simple, difficulty-driven data selection method that selects only hard examples for training. We then study how many of these examples are sufficient under this selection. Surprisingly, we find that training a 1.5B student model on only 8 selected hard examples matches the validation accuracy of the full 17K dataset baseline (53.6\% vs 53.7\%). Furthermore, we show that this extreme data efficiency generalizes robustly across diverse model architectures and scales, ranging from 1.5B to 8B parameters.
